% GENERATED FILE. Edit canonical lessons, then run npm run generate:books.
% canonical-source-hash: fnv1a64:af3217920591cd1b
% canonical-lessons: IT-C440-navigatore, IT-C440-indicazione, IT-C440-svoltare, IT-C440-proseguire, IT-C440-imboccare

\chapter{Chiedere la strada}
\label{ch:it-chiedere-la-strada}
\hlchaptermodality{440}

\begin{chapteropening}
\textbf{By the end of this chapter:} \emph{I can ask for and follow directions: the satnav, directions, turning, carrying on and taking a road.}

Complete the last lesson of chapter 440: I can ask for and follow directions: the satnav, directions, turning, carrying on and taking a road.
\end{chapteropening}

\section[il navigatore]{il navigatore — the satnav; the navigator}

\label{lesson:IT-C440-navigatore}



\begin{quote}
Before the new one: say the Italian for the street plan, then the Italian for the compass.
\end{quote}

\subsection*{You'll want to know: il navigatore}
\textbf{il navigatore} — \textquotedblleft{}the satnav; the navigator\textquotedblright{}.

\begin{quote}
Il navigatore dice di girare a sinistra.
\end{quote}

\begin{quote}
\emph{The satnav says to turn left.}
\end{quote}

In full \emph{il navigatore satellitare}; also a seafarer.

\begin{cousinweb}
\emph{navigatore} is a learned borrowing from Latin \emph{nāvigātor} 'sailor', from \emph{nāvigāre} 'to sail', built on \emph{nāvis} 'ship' (Lewis \& Short).
\end{cousinweb}

\begin{grammarlens}[title={What you've built}]
One more word for telling, explaining and describing at B1.
\end{grammarlens}

\subsection*{Your turn}
\begin{itemize}
\raggedright
  \item \emph{Say it:} \emph{il navigatore}
  \item \emph{Say it:} \emph{il navigatore}, once more
  \item \emph{Say it:} a sentence of your own with \emph{il navigatore}
  \item \emph{Recall:} say the Italian for the street plan, then the Italian for the compass, then say \emph{il navigatore} again
\end{itemize}

\begin{tcolorbox}[breakable,colback=teal!4,colframe=teal!35!black,title={Before you move on}]
What does il navigatore mean? (\textquotedblleft{}The satnav; the navigator\textquotedblright{}.) Say it once more.
\end{tcolorbox}
\section[l'indicazione]{l'indicazione — the direction (to a place); the sign}

\label{lesson:IT-C440-indicazione}



\begin{quote}
Before the new one: say the Italian for the compass, then the Italian for the satnav.
\end{quote}

\subsection*{You'll want to know: l'indicazione}
\textbf{l'indicazione} — \textquotedblleft{}the direction (to a place); the sign\textquotedblright{}.

\begin{quote}
Mi può dare un'indicazione per la stazione?
\end{quote}

\begin{quote}
\emph{Can you give me directions to the station?}
\end{quote}

Often plural: \emph{chiedere indicazioni} — to ask for directions.

\begin{cousinweb}
\emph{indicazione} is a learned borrowing from Latin \emph{indicātiō} 'a pointing out' (EtymDB), from \emph{indicāre} 'to point out', built on \emph{index} 'pointer'.
\end{cousinweb}

\begin{grammarlens}[title={What you've built}]
One more word for telling, explaining and describing at B1.
\end{grammarlens}

\subsection*{Your turn}
\begin{itemize}
\raggedright
  \item \emph{Say it:} \emph{l'indicazione}
  \item \emph{Say it:} \emph{l'indicazione}, once more
  \item \emph{Say it:} a sentence of your own with \emph{l'indicazione}
  \item \emph{Recall:} say the Italian for the compass, then the Italian for the satnav, then say \emph{l'indicazione} again
\end{itemize}

\begin{tcolorbox}[breakable,colback=teal!4,colframe=teal!35!black,title={Before you move on}]
What does l'indicazione mean? (\textquotedblleft{}The direction (to a place); the sign\textquotedblright{}.) Say it once more.
\end{tcolorbox}
\section[svoltare]{svoltare — to turn (into a street)}

\label{lesson:IT-C440-svoltare}



\begin{quote}
Before the new one: say the Italian for the satnav, then the Italian for the direction.
\end{quote}

\subsection*{You'll want to know: svoltare}
\textbf{svoltare} — \textquotedblleft{}to turn (into a street)\textquotedblright{}.

\begin{quote}
Svolta a destra dopo la farmacia.
\end{quote}

\begin{quote}
\emph{Turn right after the chemist's.}
\end{quote}

\emph{La svolta} — the turn, and a turning point.

\begin{cousinweb}
\emph{svoltare} is built from \emph{s-} and \emph{voltare} 'to turn', which is inherited from Vulgar Latin \emph{volvitāre}, built on \emph{volvere} 'to roll' (Lewis \& Short).
\end{cousinweb}

\begin{grammarlens}[title={What you've built}]
One more word for telling, explaining and describing at B1.
\end{grammarlens}

\subsection*{Your turn}
\begin{itemize}
\raggedright
  \item \emph{Say it:} \emph{svoltare}
  \item \emph{Say it:} \emph{svoltare}, once more
  \item \emph{Say it:} a sentence of your own with \emph{svoltare}
  \item \emph{Recall:} say the Italian for the satnav, then the Italian for the direction, then say \emph{svoltare} again
\end{itemize}

\begin{tcolorbox}[breakable,colback=teal!4,colframe=teal!35!black,title={Before you move on}]
What does svoltare mean? (\textquotedblleft{}To turn (into a street)\textquotedblright{}.) Say it once more.
\end{tcolorbox}
\section[proseguire]{proseguire — to carry on, to continue}

\label{lesson:IT-C440-proseguire}



\begin{quote}
Before the new one: say the Italian for the direction, then the Italian for to turn.
\end{quote}

\subsection*{You'll want to know: proseguire}
\textbf{proseguire} — \textquotedblleft{}to carry on, to continue\textquotedblright{}.

\begin{quote}
Prosegua dritto per cento metri.
\end{quote}

\begin{quote}
\emph{Carry on straight for a hundred metres.}
\end{quote}

\emph{Buon proseguimento!} — enjoy the rest of your trip.

\begin{cousinweb}
\emph{proseguire} is inherited from Latin \emph{prōsequī} 'to follow on' (EtymDB), built on \emph{sequī} 'to follow', the verb behind \emph{seguire}.
\end{cousinweb}

\begin{grammarlens}[title={What you've built}]
One more word for telling, explaining and describing at B1.
\end{grammarlens}

\subsection*{Your turn}
\begin{itemize}
\raggedright
  \item \emph{Say it:} \emph{proseguire}
  \item \emph{Say it:} \emph{proseguire}, once more
  \item \emph{Say it:} a sentence of your own with \emph{proseguire}
  \item \emph{Recall:} say the Italian for the direction, then the Italian for to turn, then say \emph{proseguire} again
\end{itemize}

\begin{tcolorbox}[breakable,colback=teal!4,colframe=teal!35!black,title={Before you move on}]
What does proseguire mean? (\textquotedblleft{}To carry on, to continue\textquotedblright{}.) Say it once more.
\end{tcolorbox}
\section[imboccare]{imboccare — to take (a road), to turn into}

\label{lesson:IT-C440-imboccare}



\begin{quote}
Before the new one: say the Italian for to turn, then the Italian for to carry on.
\end{quote}

\subsection*{You'll want to know: imboccare}
\textbf{imboccare} — \textquotedblleft{}to take (a road), to turn into\textquotedblright{}.

\begin{quote}
Imbocchiamo la prima strada a sinistra.
\end{quote}

\begin{quote}
\emph{We take the first road on the left.}
\end{quote}

Literally to put into the mouth: it also means to spoon-feed a child.

\begin{cousinweb}
\emph{imboccare} is built from \emph{in-} and \emph{bocca} 'mouth', which is inherited from Latin \emph{bucca} 'cheek' (EtymDB): going into the mouth of a street.
\end{cousinweb}

\begin{grammarlens}[title={What you've built}]
One more word for telling, explaining and describing at B1.
\end{grammarlens}

\subsection*{Your turn}
\begin{itemize}
\raggedright
  \item \emph{Say it:} \emph{imboccare}
  \item \emph{Say it:} \emph{imboccare}, once more
  \item \emph{Say it:} a sentence of your own with \emph{imboccare}
  \item \emph{Recall:} say the Italian for to turn, then the Italian for to carry on, then say \emph{imboccare} again
\end{itemize}

\begin{tcolorbox}[breakable,colback=teal!4,colframe=teal!35!black,title={Before you move on}]
What does imboccare mean? (\textquotedblleft{}To take (a road), to turn into\textquotedblright{}.) Say it once more.
\end{tcolorbox}
